\section{Complete CALM}
\label{sec:complete-calm}

We now extend CALM theory to allow a fiber to contain several outputs.  The monotonicity property we want to test as input
\ingrows from $i$ to $j\InputGe i$ must be lifted to sets: can \emph{some}
output in the fiber at $j$ \outrelbase \emph{whatever} outputs runs emit at $i$?
In addition to seal commitments on inputs, we now have a second kind of commitment: \emph{choosing}
among permitted outputs.

\subsection{Output choice by example}
\label{sec:ambiguity-story}

Figure~\ref{fig:ambiguity-frames} contrasts two specifications that permit
several outputs in one fiber.  Transaction ordering requires a choice
between incompatible serialization orders.  Traceroute observations, by
contrast, can all be \outrelpast by a larger route set.

\begin{figure}[t]
  \centering
  \small
  \setlength{\tabcolsep}{4pt}
  \begin{minipage}[t]{0.36\columnwidth}
    \raggedright
    (f) After $[A,B]$\par\smallskip
    \begin{tabular*}{\linewidth}[t]{@{}l@{\extracolsep{\fill}}l@{}}
      \toprule
      $\mathrm{in}$ & $\mathrm{out}$ \\ \midrule
      $\{A,B\}$ & \sh{$[A,B]$} \\
      \retired{$\{A,B\}$} & \retired{$[B,A]$} \\ \strikerow
      \bottomrule
    \end{tabular*}\par\smallskip
    {\footnotesize no remaining row \outrel the removed one: a commitment}
  \end{minipage}\hfill
  \begin{minipage}[t]{0.6\columnwidth}
    \raggedright
    (g) After $\{s{\to}\texttt{s1}\}$\par\smallskip
    \begin{tabular*}{\linewidth}[t]{@{}l@{\extracolsep{\fill}}l@{}}
      \toprule
      $\mathrm{in}$ & $\mathrm{out}$ \\ \midrule
      $i$ & \sh{$\{s{\to}\texttt{s1}\}$} \\
      \retired{$i$} & \retired{$\{s{\to}a\}$} \\
      \retired{$i$} & \retired{$\{s{\to}a,\ s{\to}a{\to}\texttt{s2}\}$} \\
      $i$ & \sh{$\{s{\to}\texttt{s1},\ s{\to}a\}$} \\
      $i$ & \sh{$\{s{\to}\texttt{s1},\ s{\to}a,\ s{\to}a{\to}\texttt{s2}\}$} \\
      \bottomrule
    \end{tabular*}\par\smallskip
    {\footnotesize each removed row is \outrelpast by a remaining one: an output
    retirement}
  \end{minipage}
  \caption{Two specifications with several rows at one input, each after one output
  event.  (f) Serial orders of two concurrent transactions, under the equality
  order.  (g) Route sets from a source $s$ to a service on \texttt{s1} and
  \texttt{s2}, ordered by inclusion; $s$ reaches \texttt{s1} directly and
  \texttt{s2} through hop $a$.  Legend as in Figure~\ref{fig:election-frames};
  gray alone may also mark an output retirement.}
  \label{fig:ambiguity-frames}
\end{figure}

\paragraph{Output choice and directedness.} The examples differ in what happens
to alternative outputs when one is emitted.  In panel~(f), neither
serialization order \outrel the other, so emitting one requires a
\emph{choice commitment} that removes the other while the input stays legal:
emitting $[A,B]$ chooses it over $[B,A]$.  In panel~(g), larger route sets \outrelbase both the
emitted observation and the outputs of the rows it removes.  We call such
removal an \emph{output retirement}; every removed output is
\outrelpast by a remaining one, so emitting the route observation loses no information
according to the output order and needs no
commitment: it is \emph{exposable}.  

An output event can thus remove rows in three ways: by sealing inputs, as in
Section~\ref{sec:functional-calm}; by retiring outputs, as in panel~(g); or by
choosing among outputs, as in panel~(f).  The atomic event
$\mathsf{output}(o)$ does whichever of these $o$ requires, then emits $o$;
sealing and choosing are commitments, retiring is not.

In Figure~\ref{fig:ambiguity-frames}, serialization requires commitment; traceroute avoids it.
The difference between them is the order-theoretic concept of a \emph{directed} fiber:
every finite set of outputs in a given fiber has a common \outnoun in that fiber.
Directedness at a single input is not sufficient, however;
we must also \outrelbase prior outputs as input \ingrows. \emph{Lifted Monotonicity}
combines the requirements: as input \ingrows from $i$ to $j$, some output
in the fiber at $j$ \outrel whatever outputs runs emit at $i$.  Equivalently,
outputs emitted at inputs with a common \indesc $j$ must have a common
\outnoun in the fiber at $j$: the same test as in Section~\ref{sec:election-story}, but
with many outputs per input.  Two transaction scheduling runs that emit $[A,B]$ and $[B,A]$ fail it;
two traceroute runs that emit different route sets pass it, since their union
survives in both.

\subsection{Relational specifications and states}
\label{sec:relational-specifications}

We begin with our most general definition of a specification, which admits
ambiguity of output.
\begin{definition}[Ordered relational specification]
\label{def:relational-specification}
An \emph{ordered relational specification} is
        \begin{equation}
  \Spec=(I,\InputOrd,O,\OutputOrd,S),
  \qquad S\subseteq I\times O,
  \label{eq:relational-specification}
        \end{equation}
where $(I,\InputOrd)$ and $(O,\OutputOrd)$ are partially ordered sets and $S$
is total on $I$: every input appears in at least one row.  A row $(i,o)\in S$
says that output $o$ is permitted at input $i$.  The fiber
        \begin{equation}
  \fiber{i}=\{o\in O:(i,o)\in S\}
  \label{eq:fiber}
        \end{equation}
of Section~\ref{sec:functional-specifications} may now contain several
outputs; the specification is \emph{ambiguous} when some fiber does.  A
functional specification is the special case $\fiber{i}=\{f(i)\}$.  We call
either kind simply an \emph{ordered specification}.
\end{definition}

Specification states and input events
(Sections~\ref{sec:functional-specifications} and~\ref{sec:functional-events})
apply unchanged to relational tables.

\subsection{Lifted Monotonicity}
\label{sec:lifted-monotonicity}

As input \ingrows from $i$ to $j$, two things can go wrong.  A
\emph{stranded output} is an output in the fiber at $i$ with no \outnoun in
the fiber at some \indesc $j$.  An \emph{incompatible pair} is two outputs in
the fiber at $i$ with no common \outnoun there.  We rule out
each, combine them, and restate the result in the form runs need.

First, we rule out stranded outputs.

\begin{definition}[Relational Monotonicity]
\label{def:relational-monotonicity}
An ordered relational specification satisfies \emph{Relational Monotonicity}
when every output keeps some \outnoun as input \ingrows:
        \begin{equation}
  i\InputOrd j,\ o\in\fiber{i}
  \quad\Longrightarrow\quad
  \exists u\in\fiber{j}:\ o\OutputOrd u.
  \label{eq:relational-monotonicity}
        \end{equation}
\end{definition}

For singleton fibers this is functional monotonicity,
Equation~\ref{eq:functional-monotonicity}.  It treats each output separately,
so it does not rule out incompatible pairs, which concern several outputs in
one fiber.
For a set $F\subseteq O$ of outputs, outputs $o_1,\ldots,o_k$ are
\emph{compatible in $F$} if some $u\in F$ \outrel all of them, and
\emph{incompatible in $F$} otherwise.

Second, we rule out incompatible pairs.

\begin{definition}[Directed fiber, Directed Outputs]
\label{def:fiber-directedness}
A fiber $F\subseteq O$ is \emph{directed} when every nonempty finite
$X\subseteq F$ is compatible in $F$.  An ordered relational specification has
\emph{Directed Outputs} when every fiber $\fiber{i}$ is directed.
\end{definition}

A singleton fiber is trivially directed.  As with common \indescs of inputs
(Section~\ref{sec:functional-specifications}), a common \outnoun need not be
least: directedness requires some upper bound within the fiber for each finite
set, not a least one or a fixed merge operator, so it is weaker than a
join-semilattice.  For example, the fiber
$\{\{a\},\{b\},\{a,b,c\},\{a,b,d\},\{a,b,c,d\}\}$ under inclusion is directed,
since $\{a,b,c,d\}$ \outrel every member, yet $\{a\}$ and $\{b\}$ have two
incomparable minimal common \outnouns.  Compatibility likewise does not require
comparability: two incomparable outputs may have a third output in the fiber
that \outrel both, as $\{a,b,c\}$ and $\{a,b,d\}$ do in the fiber above.  Thus
the traceroute fiber is directed: the union of any finite set of route sets in
it is again in the fiber and refines each of them. The two transaction orders in
Figure~\ref{fig:ambiguity-frames}(f), by contrast, are incompatible.

\paragraph{Lifted Monotonicity.}
To state both requirements at once, lift $\OutputOrd$ to sets.  For nonempty
$A,B\subseteq O$, define
        \begin{equation}
  A\LiftedOrd B
  \;\Longleftrightarrow\;
  \text{every nonempty finite }X\subseteq A\text{ is compatible in }B,
  \label{eq:joint-output-order}
        \end{equation}
the \emph{joint lift} of $\OutputOrd$.  We apply it with $A=\fiber{i}$ and
$B=\fiber{j}$ for $i\InputOrd j$.  When $X$ is a single output $o$, it asks
that $o$ have a \outnoun in $\fiber{j}$, which rules out stranded outputs.  When $X$
holds several outputs, it asks that they have a common \outnoun in $\fiber{j}$; with
$j=i$, this rules out incompatible pairs.

\begin{definition}[Lifted Monotonicity, monotone specification]
\label{def:lifted-monotonicity}
An ordered relational specification satisfies \emph{Lifted Monotonicity}, and
is called \emph{monotone}, when the mapping $i\mapsto\fiber{i}$ is monotone
under the joint lift:
        \begin{equation}
  i\InputOrd j
  \quad\Longrightarrow\quad
  \fiber{i}\LiftedOrd\fiber{j}.
  \label{eq:lifted-monotonicity}
        \end{equation}
\end{definition}

With singleton fibers it is functional monotonicity
(Definition~\ref{def:functional-monotonicity}).  As the reading above shows,
it implies Relational Monotonicity and Directed Outputs; conversely, for
$i\InputOrd j$ and finite nonempty $X\subseteq\fiber{i}$, Relational
Monotonicity gives each $x\in X$ a \outnoun in $\fiber{j}$, and Directed
Outputs gives those a common \outnoun in $\fiber{j}$.  So Lifted Monotonicity
holds exactly when both do.  In powerdomain terms, Relational Monotonicity
is monotonicity of $i\mapsto\fiber{i}$ under the Hoare
preorder~\cite{abramsky1994domain}, which asks only that each output in $A$
have a \outnoun in $B$; the joint lift asks that every finite set of them have
a common one.  The two agree when $B$ is directed.  Unlike the Hoare preorder,
the joint lift is not reflexive: $A\LiftedOrd A$ holds exactly when $A$ is
directed, so the transaction fiber of Figure~\ref{fig:ambiguity-frames}(f)
already fails it at $j=i$.

\paragraph{The common-\indesc condition.}
Lifted Monotonicity compares an input with its
own expansions. Runs, however, may emit outputs at incomparable inputs, say
$o_1$ at $i_1$ and $o_2$ at $i_2$, and can be reconciled only at a common
\indesc $j$ of both. The following condition asks for exactly that: a common
\outnoun of $o_1$ and $o_2$ in the fiber at $j$.

\begin{definition}[Common-\indesc condition]
\label{def:common-expansion-condition}
An ordered specification satisfies the \emph{common-\indesc condition} when,
for every finite set of rows $(i_1,o_1),\ldots,(i_k,o_k)$ of $S$ and every
common \indesc $j$ of $i_1,\ldots,i_k$, some $u\in S(j)$ \outrel every
$o_r$.
\end{definition}

For the empty set of rows the condition asks only that $S(j)$ be nonempty, which
follows from $S$ being total on $I$.

\begin{proposition}[Common-\indesc form of Lifted Monotonicity]
\label{prop:common-expansion-form}
An ordered relational specification satisfies Lifted Monotonicity if and only
if it satisfies the common-\indesc condition.
\end{proposition}

\begin{proof}[Proof sketch]
Relational Monotonicity carries each $o_r$ to some
\outnoun in the fiber at $j$, and Directed Outputs gives those a common
\outnoun there.  Full proof in Appendix~\ref{app:proofs-sec3}.
\end{proof}

\subsection{Lifted output events}
\label{sec:lifted-output-events}

For a state $(i,R)$ and $o\in R(i)$, the output event $\mathsf{output}(o)$
again keeps exactly the rows whose outputs \outrelbase $o$, ending at
$(i,R[\mathrm{out}\OutputGe o])$ as in Equation~\ref{eq:output-action}, and
emits $o$.  With several rows at one input, each removed row falls into one of
three classes.  Consider a removed row $(i',o')$.  If no row at $i'$ \outrel $o$,
the row is in $\Sealed_o$, and the event \emph{seals} $i'$.  If a kept row at
$i'$ \outrel $o'$, the row is in $\Retired_o$, and removing it is an
\emph{output retirement}.  Otherwise the row is in $\Rejected_o$: $i'$ stays
legal, but the event \emph{chooses} against $o'$.  The seal $s_o$ removes $\Sealed_o$, and the choice
$c_o$ removes $\Rejected_o$.  Both are commitments; output retirement is not.
In Figure~\ref{fig:ambiguity-frames}(f), emitting $[A,B]$ puts $[B,A]$ in
$\Rejected_o$.  In panel~(g), the two removed route sets are in $\Retired_o$.
In the functional case, $\Retired_o$ and $\Rejected_o$ are empty, so this is
the output event of Definition~\ref{def:output-event}.
Appendix~\ref{app:lifted-output-events} gives the formal definitions.

As before, $o$ is exposable when $\mathsf{output}(o)$ applies no commitment,
that is, when $\Sealed_o$ and $\Rejected_o$ are empty.  Equivalently, at every legal \indesc $j$ of $i$, $o$ and each
remaining output at $j$ have a common \outnoun at $j$
(Proposition~\ref{prop:exposability}).

\subsection{The Complete CALM theorem}
\label{sec:complete-calm-theorem}

A lifted output event ends at $(i,R[\mathrm{out}\OutputGe o])$, as in
Equation~\ref{eq:output-action}.  Specification states, input events, runs,
commitment-free runs, and consistency are defined as before, and the results
of Section~\ref{sec:executions} used output events only through that
equation.  They therefore lift unchanged: the rows remaining after a run
(Lemma~\ref{lem:run-residual}) and the survivor test
(Proposition~\ref{prop:confluence}).  With them, consistency
reduces to Lifted Monotonicity.

\begin{theorem}[Complete CALM]
\label{thm:complete-calm}
An ordered specification is consistent if and only if it satisfies Lifted
Monotonicity, equivalently the common-\indesc condition.
\end{theorem}

\begin{proof}[Proof sketch]
If the common-\indesc condition holds, a common \outnoun $u\in S(j)$
of everything emitted survives in every run.
If outputs emitted at inputs with a common \indesc $j$ are incompatible in
$S(j)$, the runs that emit them share no row at $j$ and can never be
reconciled there.  Full proof in Appendix~\ref{app:proofs-sec3}.
\end{proof}

The name CALM, for Consistency As Logical Monotonicity, tied the idea to a
logic language.  The condition that characterizes consistency is instead a
property of the specification: Consistency As Lifted Monotonicity, which for a function is
ordinary monotonicity.

The same condition governs exposability.

\begin{proposition}[Exposure boundary]
\label{prop:exposure-boundary}
An ordered specification satisfies Lifted Monotonicity if and only
if every output at every state reachable without commitment is exposable.
\end{proposition}

\begin{proof}[Proof sketch]
Under the common-\indesc condition, consistency supplies the common \outnoun
the exposure test asks for.  Conversely, the exposure test at canonical
states, one row at a time, builds the common \outnoun the condition
asks for.  Full proof in Appendix~\ref{app:proofs-sec3}.
\end{proof}

With one output per input, Lifted Monotonicity and the common-\indesc
condition are both monotonicity, so Complete CALM and
Proposition~\ref{prop:exposure-boundary} reduce to Functional CALM and
Proposition~\ref{prop:functional-exposure-boundary}.
